\section{Theory, mechanism and measurement critiques}\label{sec:mechanism}
\paragraph{Mechanisms are under-tested} \citet{kazdin2009understanding} argues that ``there is no evidence-based explanation of how or why even the most well-studied interventions produce change'', and sets standards (temporal precedence, specificity, dose--response, experiment; \citealp{kraemer2002mediators}) that most CBT mediation studies do not meet. Positive evidence exists: sudden symptom gains were preceded by cognitive change in two datasets \citep{tang1999sudden,tang2005cognitive}. It is mixed: dysfunctional-attitude change mediated symptom reduction in CBT versus interpersonal therapy, but the reverse direction fit CBT versus medication ($n=130$; \citealp{quilty2008mediational}); the CBT-side reviews call the evidence ``not yet conclusive'' \citep{driessen2010mediators} and the evidence for skill use as a mediator preliminary \citep{hundt2013skills}. Cognitive vulnerability--stress interactions predicted later depression only weakly in youth \citep{lakdawalla2007cognitive}. Defences of the cognitive model and its updates are in \citet{beck2014generic} and \citet{hofmann2013science}.

\paragraph{Is the cognitive part the active part?} In a network meta-analysis of 45 trials ($n=3{,}382$), cognitive restructuring, behavioural activation and full CBT did not differ, and all beat usual care (SMD 0.57, 0.52, 0.44) \citep{ciharova2021restructuring}. This is consistent with a behavioural reading and equally with a cognitive one; see Section~\ref{sec:ingredients} for the dismantling trials. For exposure, inhibitory learning offers an account in terms of expectancy violation rather than belief change \citep{vervliet2013extinction,weisman2018augmentation}, and a large trial found expectancy-violation-based exposure effective and optimisable \citep{pittig2021intensified}. Thought suppression can backfire \citep{wenzlaff2000suppression}, and experiential avoidance gives a third-wave counter-theory \citep{hayes1996avoidance}, though a key measure may largely track distress \citep{wolgast2014aaq} and third-wave trials were less stringent than CBT trials \citep{ost2008thirdwave}. The process-based critique of one protocol per diagnosis \citep{hayes2021processbased,hofmann2019process} comes from advocates of the alternative.

\paragraph{Common factors} Alliance ($r=.278$; \citealp{fluckiger2018alliance}), expectancy ($r=.18$; \citealp{constantino2018expectation}) and therapist effects (about 5\%; \citealp{johns2019therapist}) are replicable and modest. \citet{wampold2015common} argues they outweigh specific ingredients, and \citet{enck2019placebo} frame context effects in psychotherapy; this is contested, and it explains why trials without blinding cannot cleanly isolate CBT's specific effect.

\paragraph{Measurement} Clinician-rated scales show larger effects than self-report ($\Delta g=0.20$; \citealp{cuijpers2010selfreport}). Seven depression scales share little symptom content (mean overlap 0.36; \citealp{fried2017symptoms}); four common scales lack unidimensionality and longitudinal invariance ($n=3{,}509$; \citealp{fried2016invariance}), and sum scores hide which symptoms change \citep{fried2015sumscores}. Response is therefore partly instrument-dependent, and standardised, patient-defined outcomes are rare \citep{cuijpers2019targets,lambert2009clinical}.

\paragraph{Evidence quality} Only 11 of 115 depression trials met eight quality criteria, with $d=0.22$ versus 0.74 in the rest \citep{cuijpers2009quality}; excess significant findings and mean power of 49\% suggest small-study bias \citep{flint2014excess}; only 15\% of audited trials were prospectively registered \citep{cybulski2016registration}; and selective reporting inflated effects (SMD $-0.81$ versus $-0.54$; \citealp{miguel2021selective}). A Cochrane-method review rated all eight cognitive-therapy-versus-usual-care trials at high risk of bias \citep{jakobsen2011cttau}, though risk-of-bias tools fit psychotherapy awkwardly \citep{munder2018rob}. Fairness requires the counterweight: a multiverse analysis of 4,281 meta-analyses still gave a mean $g=0.56$, 90\% clinically relevant, though pooling all psychotherapies \citep{plessen2023multiverse}. Field-level critiques call for mechanism research and less redundant trials \citep{cuijpers2015future,tackett2019replication}.

\paragraph{Reading} R10 and the measurement and bias part of R3 are supported as limits on certainty. Nothing here shows a mechanism by which CBT specifically fails; it shows that the theory of change is under-tested and that the trial base is biased upward.
